\documentclass[12pt,a4paper]{article}

% --- Packages ---
\usepackage[utf8]{inputenc}
\usepackage{amsmath, amsfonts, amssymb}
\usepackage{graphicx}
\usepackage{booktabs} % For professional tables
\usepackage{hyperref}
\usepackage{geometry}
\usepackage{float}
\usepackage{tabularx}
\geometry{margin=1in}

% --- Metadata ---
\title{Credit Risk Modeling System: \\ Technical Research Report}
\author{Data Engineering \& Risk Management Team}
\date{\today}

\begin{document}

\maketitle

\begin{abstract}
This technical research report details the development, validation, and implementation of an end-to-end credit risk modeling system using the Lending Club dataset (2007--2018). The system encompasses a Probability of Default (PD) model via logistic regression, a two-stage Loss Given Default (LGD) model, and an Exposure at Default (EAD) model to compute Expected Loss (EL). Recent updates have aligned the framework with real-world banking and regulatory standards, including IFRS 9 three-stage Expected Credit Loss (ECL) provisioning and Basel Advanced IRB capital requirements. 
\end{abstract}

\section{Introduction}
The accurate quantification of credit risk is fundamental to the stability and profitability of financial lending institutions. This project establishes a robust pipeline capable of predicting default probabilities and calculating expected credit losses. The models operate within the regulatory context of IFRS 9 (requiring probability-weighted, forward-looking lifetime ECL) and the Basel III framework (governing risk-weighted assets and capital reserves). 

\section{Methodology}
\subsection{Data Pipeline and Preprocessing}
The dataset consists of 2.26 million loan applications. An out-of-time (OOT) split was utilized to prevent data leakage, designating 2007--2015 as the training set (831,051 loans) and 2016--2018 as the holdout testing set (538,515 loans). Weight of Evidence (WoE) encoding was applied to robustly handle non-linear relationships and missing values.

\subsection{Model Architecture}
The system employs a three-model architecture:
\begin{itemize}
    \item \textbf{Probability of Default (PD):} A logistic regression model trained with statsmodels, utilizing backward elimination for feature selection, scaled to a 300--850 credit scorecard.
    \item \textbf{Loss Given Default (LGD):} A two-stage model addressing the zero-inflated nature of recoveries (a logistic stage predicting if any recovery occurs, and a linear stage predicting the recovery magnitude).
    \item \textbf{Exposure at Default (EAD):} A linear regression predicting the Credit Conversion Factor (CCF).
\end{itemize}

\section{Results}
The pipeline was executed end-to-end on the OOT test set, yielding robust discrimination metrics but identifying necessary calibration adjustments. 

\subsection{Performance Metrics}
Table \ref{tab:results} presents the quantitative evaluation of the models on the out-of-time test dataset, highlighting the post-calibration metrics.

\begin{table}[H]
\centering
\caption{Key Performance and Portfolio Metrics (OOT Test Set)}
\label{tab:results}
\begin{tabular}{@{}ll@{}}
\toprule
\textbf{Metric} & \textbf{Value} \\ \midrule
\textbf{Discrimination (PD Model)} & \\
AUC & 0.694 \\
Gini Coefficient & 0.387 \\
Kolmogorov-Smirnov (KS) & 0.281 \\
Brier Score & 0.172 \\ \midrule
\textbf{Risk Parameters} & \\
Mean Predicted PD & 25.27\% (Calibrated) \\
Mean LGD & 92.19\% \\
Downturn LGD & 93.76\% \\ \midrule
\textbf{Portfolio Outcomes (\$3.26B Total EAD)} & \\
Expected Loss (12-month) & \$808 Million \\
IFRS 9 Lifetime ECL (Weighted) & \$1.646 Billion \\ \midrule
\textbf{Stability} & \\
Population Stability Index (Score PSI) & 0.282 (ALERT) \\ \bottomrule
\end{tabular}
\end{table}

\subsection{Visualizations}
The ROC curve demonstrates the discriminative power of the PD model on unseen data. The inclusion of these quantitative graphs ensures clarity for the stakeholders without causing text overflows.

\begin{figure}[H]
    \centering
    % Scaling width to 0.7 of text width to ensure the image does not fall out of the screen
    \includegraphics[width=0.7\textwidth]{../data/reports/L02_roc_curve.png}
    \caption{Receiver Operating Characteristic (ROC) curve for the PD model.}
    \label{fig:roc}
\end{figure}

Furthermore, the integration of IFRS 9 staging provides crucial insights into the provisioning requirements across different risk tiers.

\begin{figure}[H]
    \centering
    \includegraphics[width=0.7\textwidth]{../data/reports/L03_ifrs9_ecl_staging.png}
    \caption{IFRS 9 ECL Staging Distribution and Provisioning.}
    \label{fig:ifrs9}
\end{figure}

\section{Implementation Plan}
The rollout of the credit risk system follows a phased approach to ensure regulatory compliance and production stability. Recent milestones included intercept recalibration and scenario-weighted ECL calculations. 

Table \ref{tab:timeline} outlines the structured timeline for the next updates.

\begin{table}[H]
\centering
\caption{Phased Implementation and Rollout Milestones}
\label{tab:timeline}
\begin{tabularx}{\textwidth}{@{}lXl@{}}
\toprule
\textbf{Phase} & \textbf{Milestone Description} & \textbf{Status} \\ \midrule
\textbf{Phase 1: Core Modeling} & Finalize PD, LGD, EAD architectures and WoE encoding. & Complete \\
\textbf{Phase 2: Calibration} & Intercept log-odds shift for PD; implement Reject Inference. & Complete \\
\textbf{Phase 3: Regulatory} & Implement IFRS 9 scenario-weighted lifetime ECL logic. & Complete \\
\textbf{Phase 4: Orchestration} & Port notebook logic into Airflow DAGs and Spark jobs. & Pending \\
\textbf{Phase 5: Serving} & Deploy FastAPI \texttt{/score} endpoint with stage outputs. & In Progress \\
\textbf{Phase 6: Monitoring} & Establish Grafana dashboards for PSI and vintage analysis. & Pending \\
\textbf{Phase 7: Governance} & Finalize SR 11-7 Model Card and compliance documentation. & Pending \\ \bottomrule
\end{tabularx}
\end{table}

\section{Conclusion}
The deployed credit risk system effectively translates raw loan data into actionable risk metrics that satisfy both business optimization goals and stringent regulatory requirements (IFRS 9 and Basel III). While the model demonstrates strong discriminative power, continuous population monitoring has triggered a PSI alert (0.282), indicating significant data drift from the 2007--2015 training period to the 2016--2018 validation window. Future efforts will focus on retraining the model with more recent data and fully operationalizing the Airflow/FastAPI production infrastructure.

\end{document}
